In this section the experimental results are presented.\footnote{\label{note1}The complete collection of test VPGs can be found at:\\ https://github.com/SjefvanLoo/VariabilityParityGames/tree/master/implementation} We evaluate the performance on six sets of games:
\begin{itemize}
	\item the minepump games,
	\item the elevator games,
	\item random games of type 1 with an increasing $\lambda$,
	\item random games of type 2 with an increasing $\lambda$,
	\item random games of type 3 with an increasing $\lambda$, and
	\item random games of type 1 with an increasing number of configuration.
\end{itemize}

We present the times it took to solve a VPG. For an independent approach this means the sum of the times it takes to solve every projection of the VPG. For a collective approach this simply means the solve time for the VPG. In either case we only measure the solve time; parsing, projecting and solution printing is excluded from the evaluation.

The exact times can be found in appendix \ref{appendix:resultsexact}; in this section the results are visualized and presented in a way such that we can easily compare independent and collective approaches. We have four independent approaches:
\begin{itemize}
	\item Recursive algorithm (global),
	\item Recursive algorithm (local),
	\item fixed-point iteration (global) and
	\item fixed-point iteration (local).
\end{itemize}
For every set of problems we present four charts; for every independent approaches we present a chart where its performance is compared to one or two collective algorithms. We compare the performance of the recursive algorithm for VPGs with the performance of the original recursive algorithm and the performance of the incremental pre-solve algorithm with the performance of the fixed-point iteration algorithm. In some of the charts the solve times are divided by the independent solve times to visualize how much better or worse the collective variants perform.

The following legend holds for all charts presented in this section:
\begin{itemize}[\ ]
	\item Independent approaches:
	\begin{itemize}[\ ]
\item \raisebox{.7\height}{\begin{tikzpicture}
\path[line width=2pt,color=cyan] (19,20) edge (20,20);
\end{tikzpicture}} Recursive algorithm for parity games (global)
\item \raisebox{.7\height}{\begin{tikzpicture}
\path[line width=2pt,color=violet] (19,20) edge (20,20);
\end{tikzpicture}} Fixed-point iteration algorithm for parity games (global)
\item \raisebox{.7\height}{\begin{tikzpicture}
	\path[dashed,line width=2pt,color=cyan] (19,20) edge (20,20);
	\end{tikzpicture}} Recursive algorithm for parity games (local)
\item \raisebox{.7\height}{\begin{tikzpicture}
	\path[dashed,line width=2pt,color=violet] (19,20) edge (20,20);
	\end{tikzpicture}} Fixed-point iteration algorithm for parity games (local)
\end{itemize}
\item Collective approaches:
\begin{itemize}[\ ]
\item \raisebox{.7\height}{\begin{tikzpicture}
\path[line width=2pt,color=green] (19,20) edge (20,20);
\end{tikzpicture}} Recursive algorithm for VPGs with a symbolic representation of configurations (global)
\item \raisebox{.7\height}{\begin{tikzpicture}
\path[line width=2pt,color=red] (19,20) edge (20,20);
\end{tikzpicture}} Recursive algorithm for VPGs with an explicit representation of configurations (global)
\item \raisebox{.7\height}{\begin{tikzpicture}
\path[line width=2pt,color=orange] (19,20) edge (20,20);
\end{tikzpicture}} Incremental pre-solve algorithm (global)
\item \raisebox{.7\height}{\begin{tikzpicture}
\path[dashed,line width=2pt,color=green] (19,20) edge (20,20);
\end{tikzpicture}} Recursive algorithm for VPGs with a symbolic representation of configurations (local)
\item \raisebox{.7\height}{\begin{tikzpicture}
\path[dashed,line width=2pt,color=red] (19,20) edge (20,20);
\end{tikzpicture}} Recursive algorithm for VPGs with an explicit representation of configurations (local)
\item \raisebox{.7\height}{\begin{tikzpicture}
\path[dashed,line width=2pt,color=orange] (19,20) edge (20,20);
\end{tikzpicture}} Incremental pre-solve algorithm (local)
\end{itemize}
\end{itemize}
All the experiments are ran on a Linux x64 operating system with an Intel i5-4570 @ 3.20 GHz processor and 8GB of DDR3 RAM.
\subsection{SPL examples}
Figures \ref{fig:results_minepump} and \ref{fig:results_elevator} show the solving times (in ms) of the algorithms when applied to the SPL examples.
\begin{figure}[H]
	\centering
	\begin{multicols}{2}
		\input{"results/minepump/Zlnk product based_Zlnk fam based - explicit_Zlnk fam based - symbolic_"}\\\vfill
		\input{"results/minepump/Zlnk local product based_Zlnk fam based - explicit local_Zlnk fam based - symbolic local_"}\\
		\input{"results/minepump/Fixed-point product based_Incremental pre-solve_"}\\\vfill
		\input{"results/minepump/Fixed-point local product based_Incremental pre-solve local_"}
	\end{multicols}
	\caption{Running times on the minepump games. The x-axis shows the game numbers, these correspond with the formulas described in Table \ref{tab_minepump_formulas}. The y-axis shows, on a logarithmic scale, the number of milliseconds required to solve the VPG.}
	\label{fig:results_minepump}
\end{figure}%
For the minepump example we see that the recursive algorithm for VPGs using a symbolic representation performs particularly well; about a 3 to 18 times increase in performance compared to the independent approach. For the elevator example we also find an increase in performance for the symbolic recursive algorithm compared to the independent algorithm; about a 2 to 6 times increase. The difference is most likely because the minepump games have twice as many features as the elevator games have. Notably, for the elevator games we also find a good performance for the incremental pre-solve algorithm, which we do not find for the minepump games. Finally,  we observe that there is no clear difference between the relative performances of the global algorithms and the local algorithms.
\begin{figure}[H]
	\centering
	\begin{multicols}{2}
		\input{"results/elevator/Zlnk product based_Zlnk fam based - explicit_Zlnk fam based - symbolic_"}\\\vfill
		\input{"results/elevator/Zlnk local product based_Zlnk fam based - explicit local_Zlnk fam based - symbolic local_"}\\
		\input{"results/elevator/Fixed-point product based_Incremental pre-solve_"}\\\vfill
		\input{"results/elevator/Fixed-point local product based_Incremental pre-solve local_"}
	\end{multicols}
	\caption{Running times on the elevator games. The x-axis shows the game numbers, these correspond with the formulas described in Table \ref{tab_elevator_formulas}. The y-axis shows, on a logarithmic scale, the number of milliseconds required to solve the VPG.}
	\label{fig:results_elevator}
\end{figure}%

\subsection{Random games}
Figure \ref{fig:results_type1} shows the performance of the algorithms on type 1 games. The recursive algorithms for VPGs perform quite well, even though there are a few instances where the performance is worse than the independent approach. The symbolic variant performs quite a bit better than the explicit variant. The relative performance of the local variants of the recursive algorithms is about the same as the relative performance of the global variants. 

For the incremental pre-solve algorithm we do see a big difference between a local and global approach. The global variant performs well only for games 90 and up. The local variant, however, performs well for nearly all the games. Furthermore, even for the games where the global variant performs well relative to the independent global approach does the local variant performs even better relative to the independent local approach.
\begin{figure}[H]
	\centering
	\begin{multicols}{2}
		\input{"results/FF_randomgames/Zlnk product based_Zlnk fam based - explicit_Zlnk fam based - symbolic_"}\\
		\input{"results/FF_randomgames/Zlnk local product based_Zlnk fam based - explicit local_Zlnk fam based - symbolic local_"}\\
		\input{"results/FF_randomgames/Fixed-point product based_Incremental pre-solve_"}\\
		\input{"results/FF_randomgames/Fixed-point local product based_Incremental pre-solve local_"}
	\end{multicols}
	\caption{Running times on random games of type 1 with $\lambda = \frac{\textit{game nr}}{100}$. The x-axis shows the game numbers. The y-axis shows, on a logarithmic scale, how much faster the algorithms solved a VPG compared to the independent algorithm. Clearly the independent algorithm always has value of 1 for every VPG. If an algorithm has a value above 1 for a VPG then it performed worse than the independent algorithm; if the value is below 1 then it performed better than the independent algorithm.}
	\label{fig:results_type1}
\end{figure}%


Figure \ref{fig:results_type2} shows the performance of the algorithms on type 2 games. For the recursive algorithms we see that the explicit variant takes over from the symbolic variant. This is to be expected since these games have edge guards that are not created from features but created by picking random configurations. This decreases the performance of symbolic set operations but has no effect the performance of explicit set operations. Both variants still perform somewhat better than the independent approach. Again we do not find a significant difference between the global and local approach.

For the incremental pre-solve algorithm we find a similar result as with type 1 games. The global variant performs well only when $\lambda$ is high. The local variant performs significantly better and performs well for almost all games.
\begin{figure}[H]
	\centering
	\begin{multicols}{2}
		\input{"results/FC_randomgames/Zlnk product based_Zlnk fam based - explicit_Zlnk fam based - symbolic_"}\\
		\input{"results/FC_randomgames/Zlnk local product based_Zlnk fam based - explicit local_Zlnk fam based - symbolic local_"}\\
		\input{"results/FC_randomgames/Fixed-point product based_Incremental pre-solve_"}\\
		\input{"results/FC_randomgames/Fixed-point local product based_Incremental pre-solve local_"}
	\end{multicols}
	\caption{Running times on random games of type 2 with $\lambda = \frac{\textit{game nr}}{100}$ The x-axis shows the game numbers. The y-axis shows, on a logarithmic scale, how much faster the algorithms solved a VPG compared to the independent algorithm.}
	\label{fig:results_type2}
\end{figure}%

Figure \ref{fig:results_type3} shows the performance of the algorithms on type 3 games. For these games we see the symbolic variant of the recursive algorithm performing worse than the independent approach for almost all games. The explicit variant still performs significantly better than the symbolic variant and performs somewhat better than the independent approach. This is similar to type 2 games, which is to be expected because both types of games use configuration sets not based on features. Again we do not find a significant difference between the global and local approach.

The global incremental pre-solve algorithm performs worse than the independent approach for almost all games. Notably for these games we do not find a significant increase in relative performance when using a local variant.

Notably, the explicit recursive algorithm seems to be the only algorithm unaffected by the fact that the guard sets sizes of type 3 games vary wildly (they are distributed using a beta distribution). Maybe surprisingly, the incremental pre-solve algorithm is affected heavily by this. This is most likely because there are a lot fewer edges that admit all configurations and therefore player $\alpha$ will probably win fewer vertices in a pessimistic game for player $\alpha$.

\begin{figure}[H]
	\centering
	\begin{multicols}{2}
		\input{"results/BC_randomgames/Zlnk product based_Zlnk fam based - explicit_Zlnk fam based - symbolic_"}\\\vfill
		\input{"results/BC_randomgames/Zlnk local product based_Zlnk fam based - explicit local_Zlnk fam based - symbolic local_"}\\
		\input{"results/BC_randomgames/Fixed-point product based_Incremental pre-solve_"}\\\vfill
		\input{"results/BC_randomgames/Fixed-point local product based_Incremental pre-solve local_"}
	\end{multicols}
	\caption{Running times on random games of type 3 with $\lambda = \frac{\textit{game nr}}{100}$ The x-axis shows the game numbers. The y-axis shows, on a logarithmic scale, how much faster the algorithms solved a VPG compared to the independent algorithm.}
	\label{fig:results_type3}
\end{figure}%

\subsection{Scaling}
Figure \ref{fig:results_scalegames} shows the performance of the algorithms on type 1 games where the number of configurations increase exponentially in the x-axis of the charts. For the recursive algorithm we see that the collective approach starts outperforming the independent approach around $2^4$ configurations. As the number of configurations grow we see that the symbolic variant keeps increasing in relative performance while the explicit variants relative performance starts to flatten. This is to be expected because the performance of the explicit variant always scales linearly in the number of configurations. In the worst case the symbolic variant scales quadratically in the number of configurations, however when the sets of configurations can be represented efficiently it scales much better and in this case sublinear (since the performance of the local variant keeps increasing relative to the explicit variant).

The global incremental pre-solve algorithms does not increase notably in relative performance when the number of configurations increases. However, the relative performance of the local variant does increase in performance when the number of configurations increases. The recursion of the incremental pre-solve algorithm can be conceptualized as a tree where at every node the algorithm tries to increase the pre-solved vertices. The local variant can terminate when at some node the vertex that is being locally solved is found. In such a case the whole subtree of that node is longer computed. When the number of configurations grow then potentially the size of this subtree also grows. The fact that the local incremental pre-solve algorithm scales well in the number of configurations is most likely because the algorithm can terminate early for a relatively large set of configurations.

\begin{figure}[H]
	\centering
	\begin{multicols}{2}
		\input{"results/randomscalegames/Zlnk product based_Zlnk fam based - explicit_Zlnk fam based - symbolic_"}\\\vfill
		\input{"results/randomscalegames/Zlnk local product based_Zlnk fam based - explicit local_Zlnk fam based - symbolic local_"}\\
		\input{"results/randomscalegames/Fixed-point product based_Incremental pre-solve_"}\\\vfill
		\input{"results/randomscalegames/Fixed-point local product based_Incremental pre-solve local_"}
	\end{multicols}
	\caption{Running times on random games of type 1 with $\lambda = 0.92$ and the number of features equal to $\lfloor \textit{game nr}\rfloor$. The x-axis shows the game numbers. The y-axis shows, on a logarithmic scale, how much faster the algorithms solved a VPG compared to the independent algorithm.}
	\label{fig:results_scalegames}
\end{figure}%

\subsection{Internal metrics}
Earlier we hypothesised that the recursive algorithm for VPGs could perform well if we can attract many configurations simultaneously. For every VPG we measure the average number of configurations that were attracted simultaneously. We measure this relative to the total number of configurations in the VPG. This gives a number for every VPG. For every set of VPGs we average this number to get an average set size for every problem set. These values indicate how many configurations were attracted simultaneously. In Table \ref{tab_attracted_set_size} these values are presented for the different problems being globally solved, note that whether the sets are represented explicitly or symbolically is irrelevant. We see that this number somewhat predicts the performance of the recursive algorithms.

\begin{table}[h]
	\centering
	\begin{tabular}{|l|l|}
		\hline
		Minepump& 46\%\\ \hline
		Elevator& 51\%\\ \hline
		Type 1, scaling in $\lambda$& 61\%\\ \hline
		Type 2, scaling in $\lambda$& 55\%\\ \hline
		Type 3, scaling in $\lambda$& 22\%\\ \hline
		Type 1, scaling in \# confs& 72\%\\ \hline
	\end{tabular}
	\caption{Relative size of attracted sets}
	\label{tab_attracted_set_size}
\end{table}

The incremental pre-solve algorithm tries to outperform its independent counterpart by growing the set of pre-solved vertices. We measure how many vertices are pre-solved for the different sets of problems. For every VPG we measure the average number of pre-solved vertices for every (pessimistic) parity game solved. We measure this relative to the number of vertices in the VPG. This gives a number for every VPG. For every set of VPGs we average this number to get an average vertex size for every problem set. These values indicate how many vertices were pre-solved on average. In Table \ref{tab_mean_presolved} these values are presented for the different problems. The algorithm recurses into two branches after every two pessimistic parity games solved. The further we go down the tree the higher the number of pre-solved vertices. So the average numbers presented in the table are quite high because there are exponentially more parity games that are further down the recursion tree.

For the global variant the number presented in Table \ref{tab_mean_presolved} somewhat predicts the performance of the incremental pre-solve algorithm compared to the fixed-point iteration algorithm. For the local variant this is not the case. The local variant performs well when parts of the recursion tree are not calculated because we have terminated early. However, if the recursion tree is less deep then the numbers in the table decrease. Therefore, for the local algorithm the numbers do not predict the performance.

\begin{table}[h]
	\centering
	\begin{tabular}{|l|l|l|}
		\hline
		& Global & Local \\ \hline
		Minepump& 58\%& 29\%\\ \hline
		Elevator& 84\% & 38\%\\ \hline
		Type 1, scaling in $\lambda$& 87\%& 9\%\\ \hline
		Type 2, scaling in $\lambda$& 82\%& 9\%\\ \hline
		Type 3, scaling in $\lambda$& 57\%& 19\%\\ \hline
		Type 1, scaling in \# confs& 91\%& 12\%\\ \hline
	\end{tabular}
	\caption{Relative number of pre-solved vertices}
	\label{tab_mean_presolved}
\end{table}

\subsection{Discussion}
From the experimental results we observe that the symbolic variant of the recursive algorithms for VPGs performs well for the model verification problems. For type 1 random games it also performs well and particularly scales very well in the number of configurations. For type 2 and 3 games the sets of configurations can no longer be efficiently represented symbolically and performance drops.


After comparing independent and collective approaches we compare the performances of the algorithms overall.

First we compare the independent algorithms. Table \ref{tab_compare_independent_algs} shows for every set of VPGs how long it took each algorithm to solve all the VPGs in that set.
\begin{table}[h]
	\centering
	\begin{tabular}{|l|l|l|l|l|}
		\hline
		& Recursive &Recursive & Fixed-point & Fixed-point\\
		& global & local & global & local \\
		\hline
		Minepump& 1032 ms& 788 ms& 16713 ms& 15989 ms\\ \hline
		Elevator& 79468 ms& 65588 ms& 7393468 ms& 2108539 ms\\ \hline
		Type 1, scaling in $\lambda$& 1830 ms& 1730 ms& 53677 ms& 51506 ms\\ \hline
		Type 2, scaling in $\lambda$& 1936 ms& 1721 ms& 74710 ms& 72536 ms\\ \hline
		Type 3, scaling in $\lambda$& 1892 ms& 1680 ms& 62568 ms& 51231 ms\\ \hline
		Type 1, scaling in \# confs& 32903 ms& 29393 ms& 439199 ms& 274201 ms\\ \hline
	\end{tabular}
	\caption{Comparison of independent algorithms. The times shown are the times it took an algorithm to solve all the VPGs in a problem set independently.}
	\label{tab_compare_independent_algs}
\end{table}
We observe that the recursive algorithm performs significantly better than the fixed-point algorithm across all problems. We also see that the local variants perform somewhat better across the board.
In table \ref{tab_compare_collective_algs} we compare the performance of the collective algorithms.
We observe that the recursive symbolic variant performs the best for model-checking problems and for type 1 games. Furthermore, most likely the algorithm will scale well for models with a large number of features. The local variant of the incremental pre-solve algorithm also performs well relative to its independent counterpart. However, because the fixed-point iteration is heavily outperformed by the recursive algorithm its overall performance is worse that the recursive variants. On average, the global variant of the incremental pre-solve algorithm also outperforms its independent counterpart. However, we have seen in the comparison charts that it does so less consistently and significantly than the local variant and the symbolic recursive algorithm do.
\begin{table}[h]
	\centering
	\begin{tabular}{|l|l|l|l|l|l|l|}
		\hline
		& Recursive&Recursive &Recursive &Recursive & Incremental& Incremental\\
		& explicit &explicit & symbolic &symbolic & pre-solve &pre-solve\\
		& global & local &global  & local & global & local
		\\ \hline
		Minepump& 1019 ms& 942 ms& 148 ms& 133 ms& 5900 ms& 3223 ms\\ \hline
		Elevator& 81225 ms& 78635 ms& 25764 ms& 25602 ms& 1634659 ms& 1278387 ms\\ \hline
		Type 1, & \multirow{2}{*}{209 ms}& \multirow{2}{*}{158 ms} & \multirow{2}{*}{91 ms} & \multirow{2}{*}{86 ms} & \multirow{2}{*}{8040 ms} & \multirow{2}{*}{3801 ms}\\
		scaling in $\lambda$& & & & & &\\ \hline
		Type 2,& \multirow{2}{*}{234 ms} & \multirow{2}{*}{199 ms} & \multirow{2}{*}{2741 ms} & \multirow{2}{*}{2585 ms} & \multirow{2}{*}{67458 ms} & \multirow{2}{*}{13459 ms} \\
		scaling in $\lambda$& & & & & &\\ \hline
		Type 3,& \multirow{2}{*}{677 ms} & \multirow{2}{*}{665 ms} & \multirow{2}{*}{15891 ms} & \multirow{2}{*}{15897 ms} & \multirow{2}{*}{196328 ms} & \multirow{2}{*}{102182 ms}\\ 
		scaling in $\lambda$& & & & & &\\ \hline
		Type 1,& \multirow{3}{*}{1088 ms} & \multirow{3}{*}{1048 ms} & \multirow{3}{*}{114 ms} & \multirow{3}{*}{104 ms} & \multirow{3}{*}{53460 ms} & \multirow{3}{*}{683 ms} \\
		scaling& & & & & &\\
		in \# confs& & & & & &\\ \hline
	\end{tabular}
	\caption{Comparison of collective algorithms. The times shown are the times it took an algorithm to solve all the VPGs in a problem set collectively.}
	\label{tab_compare_collective_algs}
\end{table}

Furthermore, we observe that the explicit variant of the recursive algorithm performs decent across most games. We conclude from this that the efficiency of the symbolic algorithm does not only come from representing sets of configurations efficiently; using a collective approach even without this representation seems to be efficient. It seems that, for the games we experimented with, using the explicit recursive algorithm never significantly hurts performance but in some cases can significantly increase performance compared to the independent approach.

Earlier we hypothesized that a local-collective approach would increase performance more compared to a global-collective approach than a local-independent approach would compared to a global-independent approach. We observe that this is the case for the incremental pre-solve algorithm but not at all the case for the recursive algorithms. Furthermore, this is only the case for the random games; in the SPL games there is not a significant difference in relative performance. We conclude that local solving has the potential to increase performance, however this is not a given and depends on the algorithm and the type of VPG.